\documentclass[11pt,a4paper]{article}

\usepackage[utf8]{inputenc}
\usepackage{amsmath,amssymb,amsthm}
\usepackage{geometry}
\usepackage{hyperref}
\usepackage{booktabs}
\usepackage{xcolor}
\usepackage{enumitem}
\usepackage{url}
\usepackage{float}

\geometry{margin=2.5cm}
\hypersetup{colorlinks=true,linkcolor=blue,citecolor=blue,urlcolor=blue}

\title{Initial Design Document: RL-Driven Physical Activity Interventions\\
\large A Configurable Simulation Framework}
\author{William Dennis\\
Bristol x NUS Research Project\\
\texttt{ky23812@bristol.ac.uk}}
\date{June 2026}

\begin{document}
\maketitle
\tableofcontents

\newpage

\section{Introduction}
\label{sec:intro}

Reinforcement learning offers a principled approach to personalising
just-in-time adaptive interventions (JITAIs) for physical activity, but
training RL agents requires either expensive real-world trials or a realistic
simulation environment. Existing simulators are bespoke and tied to specific
datasets or papers, making cross-paper comparison difficult and limiting
reproducibility. We present \texttt{rl-health-interventions}, a configurable
simulation framework where the dataset schema, MDP dynamics, user behaviour
models, and agent configurations are specified entirely through YAML config
files---no source code changes required. The framework enables systematic
comparison of RL agents on PA intervention tasks using synthetic data in
Phase 1, with real-data integration planned for Phase 2.

The initial implementation delivers a rule-based user simulator with four
behavioural archetypes, a multi-timescale MDP environment (daily steps +
delayed 3-week clinical measures), and Thompson Sampling as the baseline
agent \cite{karine2024stepcountjitai,klasnja2019heartsteps}. Deep RL agents
(DQN, PPO) will be added if time permits. All
components share a common ABC + registry interface for extensibility.

\section{Literature Review}
\label{sec:lit}

Reinforcement learning has been applied to health interventions through
micro-randomized trials (MRTs), most notably HeartSteps V1 and V2, which
delivered contextually-tailored activity suggestions and demonstrated the
feasibility of RL-in-the-loop policy learning \cite{klasnja2019heartsteps,
klasnja2022heartsteps}. These trials provide the only available source of
real action-to-response mappings for calibrating user behaviour models.

Existing simulation environments for RL-driven health interventions tend to
be tied to specific datasets or tasks. StepCountJITAI \cite{karine2024stepcountjitai}
provides a simulation environment for physical activity JITAIs but is scoped
to a single dataset and action set. Health Gym \cite{gateno2023healthgym}
offers synthetic health environments but does not support configurable MDP
specifications. Open wearable benchmark datasets such as ExtraSensory
\cite{vaizman2017extrasensory} and WISDM \cite{kwapisz2011wisdm} are
valuable for pipeline testing but lack intervention response data.

The gap we address is a config-first simulation framework where the dataset
schema, MDP dynamics, user behaviour, and agent configurations are specified
entirely through config files, enabling systematic cross-dataset and
cross-policy comparisons without bespoke engineering.

\section{MDP Formalisation}
\label{sec:mdp}

This section specifies the Markov Decision Process (MDP) for the first
iteration of the simulation framework. The MDP is defined as the tuple
$(\mathcal{S}, \mathcal{A}, P, R, \gamma)$.

\textbf{State.} The state represents the current context of a user at decision epoch $t$.
States are constructed from wearable-device signals, user profile
information, and intervention history:

\[
s_t = \{ \mathbf{x}_t^{\text{wearable}}, \mathbf{x}_t^{\text{context}},
\mathbf{x}_t^{\text{history}}, \mathbf{x}^{\text{profile}} \}
\]

See Appendix~\ref{app:state} for the full list of state variables.
Variables are normalised to $[0,1]$ before being passed to the agent. The
body measure is only observed every 3 weeks (sparse reward signal).

\textbf{Actions.}
\[
\mathcal{A} = \{ a_0, a_1, a_2, a_3, a_4, a_5 \}
\]

\begin{table}[H]
\centering
\caption{Action space}
\begin{tabular}{clll}
\toprule
Action & Label & Description \\
\midrule
$a_0$ & No message & Do not intervene \\
$a_1$ & Motivational prompt & Generic encouragement \\
$a_2$ & Walking suggestion & Recommend a short walk \\
$a_3$ & Goal reminder & Remind of step goal \\
$a_4$ & Recovery suggestion & Stretch or rest \\
$a_5$ & Progress feedback & Report goal progress \\
\bottomrule
\end{tabular}
\end{table}

Exactly one action per decision epoch. The ``no message'' action is critical
for managing notification burden.

Each action carries two scalar penalties in the config schema:
\begin{itemize}[noitemsep]
    \item \texttt{reward\_penalty} --- subtracted from the activity-change reward at the current epoch.
    \item \texttt{burden\_penalty} --- added to the user's cumulative burden (fatigue model).
\end{itemize}
The ``no message'' action has zero on both axes; all other actions have
positive values. Penalties are configurable per action, allowing researchers
to model different notification intensities (e.g. a walking suggestion may
have higher burden cost than a motivational prompt). Whether these two
penalties should be merged into a single mechanism is an open question
--- see Decision Log \S\ref{sec:decisions}.

\textbf{Transition.} The transition function models the probability over next states:

\[
P(s_{t+1} \mid s_t, a_t): \mathcal{S} \times \mathcal{A} \times \mathcal{S}
\rightarrow [0,1]
\]

In Phase 1 (rule-based), transitions follow a behavioural response model:

\[
\Delta \text{steps}_{t+1} = f(s_t, a_t, \theta_{\text{user}})
\]

where $\theta_{\text{user}}$ parameterises one of four behavioural archetypes:

\begin{itemize}[noitemsep]
    \item \textbf{Goal-driven:} Responds to reminders and feedback.
    Proportional to goal progress.
    \item \textbf{Social responder:} Responds to motivational prompts.
    Models social desirability bias.
    \item \textbf{Resistant:} Low response. Fatigue accumulates quickly.
    \item \textbf{Stable maintainer:} Already active. Low marginal gain.
\end{itemize}

Burden evolves as:

\[
\text{burden}_{t+1} = \max(0,\; \text{burden}_t - \delta \cdot \mathbb{1}[a_t = a_0] + \text{burden\ penalty}_{a_t})
\]

where $a_0$ is the ``no message'' action, $\delta > 0$ is the configurable
burden decay rate, and $\mathbb{1}[\cdot]$ is the indicator function.
Burden decays only when no intervention is sent, allowing user recovery.
When $\text{burden} > B_{\max}$, response probability decays linearly.
This is a simple linear-accumulator model. The closest published precedent
is StepCountJITAI's bounded habituation variable \cite{karine2024stepcountjitai},
which uses a $[0,1]$ increment/decrement formulation. The most appropriate
burden model for PA interventions is an open question
--- see Decision Log \S\ref{sec:decisions}.

\textbf{Reward.} The reward has an immediate and a delayed component.

Immediate reward:
\[
R_t^{\text{immediate}} = \alpha \cdot \Delta\text{activity}_t
- \beta \cdot \text{reward\_penalty}_{a_t}
+ \lambda \cdot \text{goal\_progress}_t
\]
where $\Delta\text{activity}_t$ is a configurable activity metric (step
count by default; active minutes, METs, or a composite can be substituted
via config). Whether steps is the most appropriate primary metric is an
open question --- see Decision Log \S\ref{sec:decisions}.

Delayed reward (body measures) arrives every 3 weeks ($t = 21k$, i.e.\ every 21 days):

\[
R_{21k}^{\text{delayed}} = \eta \cdot \text{BM\_improvement}_{21k}
\]

Whether this sparse formulation is optimal, or whether a decaying reward
(e.g.\ $\eta \cdot \text{BM\_improvement} \cdot \gamma^{21k-t}$ between
body measures) would improve learning, is an open question
--- see Decision Log \S\ref{sec:decisions}.

Total reward:

\[
R_t =
\begin{cases}
R_t^{\text{immediate}} + R_t^{\text{delayed}} & t \equiv 0 \pmod{21} \\
R_t^{\text{immediate}} & \text{otherwise}
\end{cases}
\]

The compound reward (immediate + delayed) is a central research question.
Weights $(\alpha, \beta, \lambda, \eta)$ are configurable experiment parameters.

\textbf{Discount factor.} $\gamma \in [0.9, 0.99]$. A high discount factor
encourages long-term behaviour change over short-term step increases.
Appropriate for health habit formation. The optimal range for PA interventions
is an open question --- see Decision Log \S\ref{sec:decisions}.

\textbf{Decision epoch.} Daily decisions. The framework provides a base library
of policies; per-user policy refinement operates within these constraints.
Daily decisions align with typical intervention cadence, the 3-week body
measure delay (21 epochs), and practical deployment constraints. Whether
hourly decisions would be more appropriate is an open question
--- see Decision Log \S\ref{sec:decisions}.

\section{Framework Design}
\label{sec:design}

The framework is organised into six modules that form a pipeline from
configuration to results. All extension points use an ABC + registry pattern:
implement an abstract class, register it with a key, and reference that key in
config --- no source code changes required to add new datasets, agents, or
behaviour models.

A typical experiment starts with a YAML file that selects the dataset,
MDP parameters, user profiles, and agent. The framework wires everything
from config --- adding a new dataset or agent means writing a new config
(and optionally a new ABC subclass), not touching framework code.

\begin{table}[H]
\centering
\caption{Software modules and their dependencies}
\begin{tabular}{lll}
\toprule
Module & Role & Depends on \\
\midrule
Config & YAML $\rightarrow$ Pydantic schemas, 3-layer validation & --- \\
Data & Polars lazy reader, FeaturePipeline, StateView & Config \\
MDP & Environment step/reset, Transition ABC, Reward ABC & Data \\
UserSim & UserProfile (4 archetypes), ResponseModel ABC & MDP \\
Agent & Agent ABC, ThompsonSampling (+ optional DQN) & MDP \\
Experiment & Factory wiring, CLI, results output & All above \\
\bottomrule
\end{tabular}
\end{table}

The ExperimentFactory reads a composed experiment config, instantiates
each component in dependency order, and runs the trial loop. Results are
written to CSV with a YAML sidecar for reproducibility. The ABC + registry
pattern means swapping any component requires only a new subclass and a
config key --- no framework modifications. See
\texttt{docs/code/codebase\_plan.md} and \texttt{docs/code/code\_design.md}
for the full specification and \texttt{README.md} for a milestone/issue
dependency map.

\section{Data Sources}
\label{sec:data}

This project targets two tiers of data. \textbf{Phase 1} uses synthetic data
generators parameterised from published population statistics (e.g., NHANES
step counts, All of Us sleep distributions). This allows pipeline development and
agent benchmarking while real data access is being arranged. Unified data
loaders for all open datasets are provided by the data layer (see
\texttt{src/rl\_health\_interventions/data/}).

\textbf{Phase 2} targets three real datasets for validation:
HeartSteps V1/V2 \cite{klasnja2019heartsteps,klasnja2022heartsteps}
provide MRT data with intervention response mappings for calibrating the
user simulator. The All of Us Fitbit Dataset \cite{patten2026allofus}
offers 59,000+ participants with wearable data and EHR linkage. The UK
Biobank Accelerometer Dataset \cite{williamson2024ukbiobank} provides
700,000+ person-days of raw accelerometer data.

Detailed feasibility studies are in \texttt{docs/sources/data\_sources.md}
and \texttt{docs/sources/additional\_data\_sources.md}.

\section{Decision Log}
\label{sec:decisions}
\begin{table}[H]
\centering
\begin{tabular}{p{0.30\textwidth}p{0.55\textwidth}c}
\toprule
Question & Decision & Status \\
\midrule
What config format? & YAML & Confirmed \\
What interface? & CLI + config files; web UI = stretch & Confirmed \\
Which language and package manager? & Python 3.11 + \texttt{uv} & Confirmed \\
What data for Phase 1? & Synthetic data while awaiting real data access & Open \\
Which agent first? & Thompson Sampling & Confirmed \\
How often are decisions made? & Daily epochs & Confirmed \\
What reward structure? & Multi-timescale: immediate + 3-week delayed & Confirmed \\
Which user archetypes? & Goal-driven, social, resistant, stable maintainer & Confirmed \\
How are components wired? & ABC + registry pattern & Confirmed \\
Which dependencies? & Minimal: no Gymnasium, no SB3 & Confirmed \\
How to distribute? & Both: repo-based now, package if needed later & Confirmed \\
Is MDP formalisation reasonable? & Awaiting supervisor approval & Open \\
Which datasets for Phase 2? & HeartSteps first, then All of Us, UK Biobank & Open \\
What format for experiment output? & CSV, terminal table, JSON, or summary plots & Open \\
Metrics for success? & Output metrics to evaluate agent performance (regret, reward, adherence) & Open \\
What action penalties for which archetypes? & reward\_penalty and burden\_penalty per action, per archetype & Open \\
What discount factor range is optimal? & $\gamma \in [0.9, 0.99]$ to be explored empirically & Open \\
Decision epoch frequency? & Daily (base library of policies; per-user refinement within constraints). Hourly is an alternative. & Open \\
Merge reward penalty and burden penalty? & Currently separate (different roles: direct reward vs. state-accumulator). Merging is possible if simpler. & Open \\
Which burden model is most appropriate? & Linear accumulator $\max(0, burden_t - \delta \cdot \mathbb{1}[a_t=a_0] + penalty)$ with threshold decay. Configurable decay rate $\delta$. StepCountJITAI uses bounded $[0,1]$ habituation. & Open \\
What activity metric for the immediate reward? & Steps (default). Active minutes, METs, or composites are configurable alternatives. & Open \\
Sparse vs decaying delayed reward? & 3-week sparse (current). Decaying formulation ($\eta \cdot BM \cdot \gamma^{21k-t}$) may improve credit assignment. & Open \\
\bottomrule
\end{tabular}
\end{table}

\section{Evaluation Plan}
\label{sec:evaluation}

\textbf{Baselines.} We compare Thompson Sampling against three baselines:
(1)~random policy (uniform action selection), (2)~fixed policy (always send
motivational prompt), and (3)~rule-based policy (intervene when steps fall
below threshold). These baselines span the spectrum from no intelligence to
domain-expert heuristics.

\textbf{Metrics.} Primary: cumulative regret (lower is better), average
reward per episode. Secondary: adherence rate (fraction of epochs with
intervention acceptance), sustained behaviour change (percentage of users
maintaining $>$baseline activity for 3+ months), intervention burden
(average notifications per user per week).

\textbf{Statistical analysis.} We use bootstrap confidence intervals (1000
resamples) for pairwise comparisons between agents. Effect sizes reported
as Cohen's $d$ with 95\% CIs. Multiple comparison correction via
Bonferroni--Holm procedure. Power analysis: minimum detectable effect size
of 10\% improvement over best baseline with 80\% power at $\alpha=0.05$.

\textbf{Experimental protocol.} Each configuration runs 100 episodes with
100 synthetic users $\times$ 90 days. Seeds fixed for reproducibility.
Ablation studies isolate reward components ($\alpha$, $\beta$, $\lambda$,
$\eta$) and transition model effects. Results reported with mean $\pm$
standard deviation across seeds.

\section{Safety \& Ethics Considerations}
\label{sec:safety}

\textbf{Safety constraints.} The framework enforces hard safety limits:
(1)~maximum intervention frequency (configurable, default 3 per day),
(2)~minimum recovery period between interventions (configurable, default
2 hours), (3)~burden threshold $B_{\max}$ that blocks interventions when
exceeded. Safety violations logged with timestamps for post-hoc analysis.
These constraints prevent notification fatigue and ensure user recovery.

\textbf{Ethical approval.} Phase~1 uses synthetic data only and does not
require IRB approval. Phase~2 targets publicly available datasets
(HeartSteps V1/V2) and datasets with existing data use agreements
(All of Us, UK Biobank). We will document IRB/ethics approval status for
each dataset before access. No new human-subjects data collection is
planned.

\textbf{Privacy \& data governance.} All Phase~2 datasets are de-identified.
We comply with dataset-specific data use agreements. For GDPR/HIPAA: no
PHI/PII stored locally, data access logged, right to erasure respected via
dataset provider protocols. Data provenance tracked via config snapshots
and version control.

\textbf{Bias \& fairness.} The 4 user archetypes may not capture demographic
variation. We will: (1)~document archetype parameter sources, (2)~test
policy performance across archetype distributions, (3)~report disparate
impact metrics if demographic data available. Health equity considerations
documented in design review.

\section{Comprehensive Audit Summary}
\label{sec:audit}

A comprehensive 6-phase audit (June 2026) assessed the framework's readiness
for implementation and Nature Methods publication. The audit covered design
analysis, documentation alignment, roadmap creation, issue generation, and
6-domain code quality assessment.

\textbf{Key findings.} Overall repo health: 3.3/10. The framework has
excellent architectural scaffolding (ABC + registry pattern, config-first
design) but zero functional implementation. 8 of 17 source files are pure
stubs returning default values. All ABC interfaces use \texttt{Any} type
annotations, defeating type safety. Documentation is 65\% aligned with code,
but subphase status fields are systematically stale.

\textbf{Issue breakdown.} 43 issues generated across 4 domains:
(1)~10 milestone issues (M-01 through M-10) covering config schemas,
StateView/Environment, transition/reward models, user simulation, Thompson
Sampling, experiment runner, synthetic data, evaluation framework,
documentation, and safety/ethics;
(2)~13 risk issues including 5 RL-specific risks (reward hacking,
distribution shift, safety violations, data leakage, clinical validity
failures);
(3)~8 documentation debt issues (stale status fields, missing interfaces,
misleading README, missing dataset loaders);
(4)~12 stub implementation issues (7 existing stubs requiring implementation,
5 missing modules requiring creation).

\textbf{Critical blockers for publication.} (1)~No working experiment
runner---framework cannot execute experiments. (2)~No evaluation
methodology---no baselines, metrics, or statistical plan. (3)~No safety
constraints---health interventions require explicit safety guarantees.
(4)~No ethics/privacy documentation---Phase~2 requires real health data.
(5)~All ABC interfaces use \texttt{Any} types---type checking provides no
safety. (6)~No CI/CD pipeline---quality cannot be maintained. (7)~No
example experiment configs---users cannot learn framework usage.

\textbf{Roadmap.} 10 milestones across 4 sprints (12--16 weeks to Nature
submission). Sprint~1 (weeks 1--3): config schemas, StateView/Environment,
synthetic data, CI/CD. Sprint~2 (weeks 4--6): transition/reward models,
user simulation, Thompson Sampling. Sprint~3 (weeks 7--9): experiment
runner, evaluation framework, documentation. Sprint~4 (weeks 10--12):
safety constraints, ethics documentation, clinical validity. Nature Methods
readiness currently 14\% (2/14 requirements met); target 100\% by Sprint~4.

\textbf{Recommended first action.} Implement StateView dataclass and
Environment class with step/reset API (Milestone M-02). This unblocks all
RL experimentation and is on the critical path. Estimated effort: 12 hours.

Full audit reports available in \texttt{reports/} directory and
\texttt{AUDIT\_MASTER.md} at repository root.

\bibliographystyle{unsrt}
\bibliography{references}

\appendix
\section{State Variables}
\label{app:state}

\begin{table}[H]
\centering
\caption{State variables}
\begin{tabular}{lll}
\toprule
Variable & Type & Description \\
\midrule
$\text{steps}_t$ & $\mathbb{R}_{\geq 0}$ & Step count, past 24h \\
$\text{hr}_t$ & $\mathbb{R}_{\geq 0}$ & Heart rate (or resting) \\
$\text{sleep\_hours}_t$ & $\mathbb{R}_{\geq 0}$ & Previous night sleep \\
$\text{sedentary\_min}_t$ & $\mathbb{R}_{\geq 0}$ & Sedentary minutes today \\
$\text{time\_of\_day}_t$ & $\{0,\dots,23\}$ & Current hour \\
$\text{day\_of\_week}_t$ & $\{0,\dots,6\}$ & Day of week \\
$\text{goal\_progress}_t$ & $[0,1]$ & Fraction of daily goal met \\
$\text{burden}_t$ & $\mathbb{R}_{\geq 0}$ & Cumulative burden (decays when no interventions) \\
$a_{t-1}$ & $\mathcal{A}$ & Previous action \\
$\text{response}_{t-1}$ & $\{0,1\}$ & Previous response \\
$\text{body\_measure}_k$ & $\mathbb{R}$ & Clinical measure (q3wks) \\
$\text{age}$ & $\mathbb{Z}_{>0}$ & Static profile \\
$\text{gender}$ & $\{0,1\}$ & Static profile \\
$\text{baseline\_activity}$ & $\{0,1,2\}$ & Low / medium / high \\
\bottomrule
\end{tabular}
\end{table}

\end{document}
